\section{A Certified, Attributed Evaluation Protocol}\label{sec:protocol}

Disprovers can only be compared if every tool is judged by the same semantics, the same admissibility rules and the same budget, and if their output is judged for meaning as well as counted. Amendments and post-hoc deviations are listed with timestamps in the artifact repository and summarised in Section~\ref{sec:threats}.

\textbf{Independent certifier.} A module of about 450 lines, plus a 300-line profile checker, that imports no generator code. Given a witness (rows per table), it executes the \emph{original} DDL on a fresh in-memory SQLite database, inserts the rows with parameter binding, checks foreign keys at value level (some BIRD foreign keys reference non-unique columns, so SQLite's enforcement cannot be used), runs the \emph{original} gold and predicted strings without rewriting, and evaluates the contracts and tiers of Section~\ref{sec:problem}. Every witness of every tool, including the counterexamples released by SpotIt and SpotIt+, passes through it. Lexical formats for $\mathcal A_{1f}$ come from the value samples that the SpotIt+ authors mined from the official BIRD databases (123 text columns).

\textbf{Populations.} The primary population is the EX-correct (prediction, gold) pairs; the secondary population, used for recall, is the EX-incorrect pairs whose prediction executes. Predictions identical to the gold query up to whitespace cannot differ and count as unresolved for every tool. Results on distinct (question, gold, prediction) triples are mapped back to every system that produced the triple.

\textbf{Budgets and acceptance.} Every self-run tool receives the same wall-clock budget per distinct pair on one core, and every tool's acceptance test is the certifier at the evaluated tier: a tool that generates a non-admissible distinguishing database keeps searching until the budget expires. For SpotIt and SpotIt+ we use the authors' released per-question results (bounds up to five, median 5.7\,s, 95th percentile about 20 minutes per pair). We report them in two readings: the \emph{author-claimed} count of non-equivalent pairs and the \emph{re-certified} count of released counterexamples that our certifier accepts. Hypothesis tests use the claimed counts, the reading most favourable to the baselines.

\textbf{Attribution.}\label{sec:attribution} For each tool we draw a random sample of 100 distinct certified pairs at its primary tier ($\mathcal A_0$ for the released SpotIt/SpotIt+ counterexamples, which rarely reach $\mathcal A_{1f}$; $\mathcal A_{1f}$ or $\mathcal A_R$ otherwise) and classify the witness's difference into the classes of Section~\ref{sec:meaningful} (prediction errors include pairs where both queries are wrong). Samples of all tools of one stage are pooled and shuffled, so annotators do not see which tool produced a witness; blinding is imperfect, since synthetic values (\texttt{'x1'}, \texttt{-1}) are recognisable and annotators could check whether a kind of data occurs by running read-only SQL against $D_0$. Every item is annotated twice independently under the same written instructions, and disagreements are adjudicated by a third pass. Annotation was performed by an AI model (Claude, Anthropic; the model that also implemented the tools); the first author checked two samples by hand (Section~\ref{sec:threats}). The estimated number of meaningful witnesses of a tool is its certified count times its sampled share. A sample of 100 pairs covers \V{attrQMin}--\V{attrQMax} distinct questions, so we report intervals from a bootstrap over questions.

\textbf{Held-out split and three stages.} The 1,534 BIRD development questions are split once by a salted hash: three schemas are held out entirely, and 30\,\% of the questions of the other eight schemas, stratified by schema and difficulty, form the development set; questions inspected in earlier phases of the project were forced into development, and eleven such questions of held-out schemas were excluded. In the \emph{first stage}, methods were developed on development questions only; the protocol and hypotheses H1--H4 were fixed, the code was frozen (commit and SHA-256 hashes recorded), and the test questions were run once. The data-profile tier and CaliberGraph-R form the \emph{second stage}. They were conceived after we had seen the attribution of first-stage \emph{test} witnesses, so the idea of the tier and its constraint classes was informed by test data (several classes, e.g., symmetric bonds and lookup tables, bind mainly on held-out schemas). Their implementation and thresholds were developed on development questions only, hypotheses H5--H7 were fixed, the code was frozen a second time, and the test questions were run once. Section~\ref{sec:rresults} therefore also reports second-stage results separately for held-out and other schemas. A \emph{third stage}, pre-registered after both test runs and before any of its data were seen, tests the validity of the attribution and replicates the second stage on Spider, which no design decision had seen, and on thirteen unseen BIRD systems.

\textbf{Pre-registered hypotheses.} \emph{First stage:} \textbf{H1} CaliberGraph certifies more EX-correct pairs at $\mathcal A_{1f}$ than SpotIt claims; \textbf{H2} CaliberGraph+SC (CaliberGraph with the SpotIt+ constraints) certifies more at $\mathcal A_2$ than SpotIt+ claims; \textbf{H3} CaliberGraph's recall on known positives exceeds those of SpotIt (claims) and ParSEval; \textbf{H4} removing any one of anchoring, isolation, join-graph completion or the SQLite value domain significantly lowers the count. \emph{Second stage:} \textbf{H5} CaliberGraph-R certifies more pairs at $\mathcal A_R$ than random fuzzing-R and ParSEval-R; \textbf{H6} its estimated number of meaningful witnesses exceeds theirs; \textbf{H7} its share of implausible witnesses is lower than CaliberGraph's at $\mathcal A_{1f}$. \emph{Third stage:} \textbf{A} blind human annotators agree with each other and with the AI ($\kappa_2\ge0.6$) and find more implausible witnesses in the first stage than in the second (repeated in a separately pre-registered round~\textbf{R2} with facts of $D_0$); \textbf{B} a second model agrees with the AI and shows the same stage difference; \textbf{C} our issue list predicts the corrections of Jin et al.~\cite{jin2026annotation}; \textbf{S1}--\textbf{S3} repeat H5--H7 on Spider; \textbf{E1} repeats H5 on unseen systems. Two further pre-registrations followed reviews: \textbf{R3}, a third human round with three annotators and 150 witnesses (agreement $\ge0.6$ among them and with the AI, the stage gap, and a different prediction-error share from the AI's), and \textbf{SP}, a local re-run of SpotIt under our budget (CaliberGraph certifies more at $\mathcal A_{1f}$ and $\mathcal A_0$).

\textbf{Statistics.} Paired outcomes (witness or not) are compared with exact McNemar tests, and confidence intervals of count differences come from a bootstrap that resamples \emph{questions}, since the ten systems answer the same questions; because McNemar tests treat pairs as independent, the question-level intervals are the primary evidence and we print $p$-values below $10^{-10}$ as such. H1--H3 form a Holm-corrected family (H3 as an intersection of two tests), as do the two comparisons of H5; for H4 we add Holm-adjusted $p$-values over the four components post hoc. H6 was pre-registered with a parametric bootstrap from the attribution samples; we also report a bootstrap over questions. H7 is a one-sided two-proportion test.
